\documentclass[12pt,a4paper]{article}

% ── Packages ──────────────────────────────────────────────────────────────────
\usepackage[top=2.5cm,bottom=2.5cm,left=3cm,right=2.5cm]{geometry}
\usepackage[T1]{fontenc}
\usepackage[utf8]{inputenc}
\usepackage{microtype}
\usepackage{setspace}
\usepackage{titlesec}
\usepackage{booktabs}
\usepackage{tabularx}
\usepackage{array}
\usepackage{graphicx}
\usepackage{float}
\usepackage{caption}
\usepackage{amsmath}
\usepackage{siunitx}
\usepackage{xcolor}
\usepackage{hyperref}
\usepackage{natbib}
\usepackage{enumitem}
\usepackage{fancyhdr}
\usepackage{soul}

% ── Hyperref setup ─────────────────────────────────────────────────────────────
\hypersetup{
  colorlinks=true,
  linkcolor=black,
  citecolor=blue!60!black,
  urlcolor=blue!60!black,
  pdftitle={LCA of Small-Scale Tofu Production in Semarang City},
  pdfauthor={Abdul Jabbar, Putri Alifa Kholil, Desiana Fitri Awati}
}

% ── Typography ─────────────────────────────────────────────────────────────────
\onehalfspacing
\setlength{\parindent}{0pt}
\setlength{\parskip}{7pt}

% ── Section formatting ─────────────────────────────────────────────────────────
\titleformat{\section}{\large\bfseries}{\thesection.}{0.5em}{}
\titleformat{\subsection}{\normalsize\bfseries}{\thesubsection}{0.5em}{}
\titleformat{\subsubsection}{\normalsize\itshape}{\thesubsubsection}{0.5em}{}
\titlespacing{\section}{0pt}{14pt}{6pt}
\titlespacing{\subsection}{0pt}{10pt}{4pt}
\titlespacing{\subsubsection}{0pt}{8pt}{3pt}

% ── Header / Footer ────────────────────────────────────────────────────────────
\pagestyle{fancy}
\fancyhf{}
\fancyhead[L]{\small\textit{Jabbar et al. / LCA of Tofu Production in Semarang City}}
\fancyhead[R]{\small\thepage}
\renewcommand{\headrulewidth}{0.4pt}

% ── Caption style ──────────────────────────────────────────────────────────────
\captionsetup{font=small, labelfont=bf, labelsep=period}
\sisetup{output-exponent-marker=\ensuremath{\mathrm{E}}}

% ── Colours ───────────────────────────────────────────────────────────────────
\definecolor{headbg}{HTML}{e8f4f8}
\definecolor{ruleblue}{HTML}{457b9d}
\definecolor{warnamber}{HTML}{b45309}

% ══════════════════════════════════════════════════════════════════════════════
\begin{document}
% ══════════════════════════════════════════════════════════════════════════════


% ── Title block ───────────────────────────────────────────────────────────────
\begin{center}
  {\LARGE\bfseries Life Cycle Assessment of Small-Scale Tofu Production\\[4pt]
  in Semarang City: Hotspot Analysis Using CML-IA and AWARE Methods}

  \vspace{1em}

  {\large Abdul Jabbar$^{1}$,\; Putri Alifa Kholil$^{1,*}$,\; Desiana Fitri Awati$^{1}$}

  \vspace{0.4em}
  {\small $^{1}$Study Program of Environmental Science, Faculty of Mathematics and Natural Sciences,\\
  Universitas Negeri Semarang, Semarang, Indonesia}

  \vspace{0.3em}
  {\small $^{*}$Corresponding author: \href{mailto:putrialifa@mail.unnes.ac.id}{putrialifa@mail.unnes.ac.id}}
\end{center}

\vspace{0.6em}
\noindent\textcolor{ruleblue}{\rule{\linewidth}{1.2pt}}

% ── Abstract ──────────────────────────────────────────────────────────────────
\vspace{0.5em}
{\small
\noindent\textbf{Abstract.}\quad
The tofu industry is one of the most important food-processing sectors in Indonesia, given its high
consumption rate and notable nutritional value. However, the environmental impacts of each production
stage have not been comprehensively identified. Life Cycle Assessment (LCA) with a gate-to-gate system
boundary was applied to evaluate the environmental impacts of tofu production at a small and
medium-sized enterprise (SME) in Semarang City. The LCA was conducted using SimaPro~10.3 with the
CML-IA Baseline method, covering five impact categories: Global Warming Potential (GWP), Ozone Layer
Depletion (ODP), Fresh Water Aquatic Ecotoxicity (FAET), Acidification, and Eutrophication. Water use
was additionally assessed using the AWARE method. The results demonstrate that the washing and boiling
stages are the primary contributors to environmental impacts across all assessed categories. Normalization
results indicate that Fresh Water Aquatic Ecotoxicity is the most dominant impact category
(\num{1.33e3}~kg 1,4-DB~eq), followed by Global Warming Potential (\num{5.28e3}~kg CO$_2$~eq). The
principal environmental hotspots are the organic wastewater generated during the washing, soaking, and
coagulation stages, as well as firewood combustion during boiling. Recommended improvement strategies
include the implementation of wastewater treatment systems, water recycling practices, and fuel
substitution.

\vspace{4pt}
\noindent\textbf{Keywords:} Life Cycle Assessment; tofu production; CML-IA Baseline; AWARE;
environmental hotspot; Semarang
}

\noindent\textcolor{ruleblue}{\rule{\linewidth}{1.2pt}}
\vspace{0.4em}

% ══════════════════════════════════════════════════════════════════════════════
\section{Introduction}
% ══════════════════════════════════════════════════════════════════════════════

Industrial activity is one of the primary drivers of global economic growth \citep{Ghisellini2016}.
The tofu industry represents a significant food-processing sector in Indonesia owing to its high
consumption rate and notable nutritional value. According to Statistics Indonesia (BPS), average per
capita tofu consumption reached approximately 0.163~kg per week in 2024, indicating that tofu remains
a staple food in many Indonesian households \citep{Dewi2024}. This high demand has driven increased
production activity, resulting in greater resource and energy consumption, as well as the generation of
waste and emissions that may cause environmental pollution if not properly managed \citep{Rosyidah2020}.

Such impacts are particularly evident in river water quality, as the discharge of industrial wastewater
can cause Biochemical Oxygen Demand (BOD) and Chemical Oxygen Demand (COD) levels to exceed regulatory
standards. In the case of the tofu industry in Mojokerto, liquid waste was reported to contain BOD
levels exceeding 2{,}000~mg/L and COD levels exceeding 5{,}000~mg/L, far above the standard thresholds
of approximately 50 to 100~mg/L \citep{Sjafruddin2022}. Furthermore, the environmental impacts
associated with each stage of tofu production have not yet been comprehensively identified. A systematic
approach such as Life Cycle Assessment (LCA) is therefore required to evaluate the full environmental
burden of the tofu production process \citep{Islami2025}.

Life Cycle Assessment (LCA) is a standardized method used to systematically evaluate the environmental
impacts of a product or process based on all input and output flows throughout its life cycle
\citep{Hauschild2017}. This method is capable of identifying production stages that constitute
environmental hotspots, thereby serving as the scientific basis for developing targeted environmental
improvement strategies \citep{Chitaka2023}. Previous studies have applied LCA to assess the
environmental impacts of the tofu industry; however, these studies predominantly focused on climate
change impact categories and thus failed to comprehensively evaluate other environmental concerns, such
as water pollution loads, water scarcity, and solid waste generation across all production stages
\citep{Nugroho2024, Wardana2024}.

Therefore, this study aims to analyze the environmental impacts of tofu production in Semarang City
using an LCA approach with a gate-to-gate system boundary, encompassing all stages from washing to
packaging, in order to identify environmental hotspots and formulate more sustainable improvement
strategies.



% ══════════════════════════════════════════════════════════════════════════════
\section{Methods}
% ══════════════════════════════════════════════════════════════════════════════

This research utilized a quantitative methodology based on Life Cycle Assessment (LCA) to assess the
environmental impacts of tofu production at an SME in Semarang City. The evaluation was performed in
compliance with \textbf{SNI ISO~14040:2016} and \textbf{SNI ISO~14044:2016}, comprising four sequential
phases: (1)~goal and scope definition, (2)~life cycle inventory analysis, (3)~life cycle impact
assessment, and (4)~interpretation. Primary data were obtained through direct field measurements,
structured interviews with the facility operator, and on-site documentation during a complete production
batch. Secondary data were sourced from peer-reviewed scientific literature and the Ecoinvent~v3
database within SimaPro~10.3 \citep{Hauschild2017}.

\subsection{Goal and Scope Definition}

This study aims to assess the environmental impacts of tofu production at the facility level, focusing
on water use, energy consumption, and waste generation at each production stage. The \textbf{functional
unit} was established as \textbf{1~kg of tofu produced}, consistent with functional units used in
comparable LCA studies of tofu and soy-based food processing. All inventory data were initially
collected on a per-batch basis representing one complete production day, with 1{,}400~kg of soybeans as
raw material input yielding 3{,}681~kg of tofu.

The \textbf{system boundary} adopts a gate-to-gate approach, encompassing all unit processes from the
initial washing of soybeans to final molding and packaging (Figure~\ref{fig:sysboundary}). Upstream
operations, including soybean farming and transportation, as well as downstream activities such as
distribution, consumption, and end-of-life management, were excluded from the system boundary. This
delineation is consistent with prior gate-to-gate LCA investigations of similar Indonesian tofu SMEs
\citep{Nugroho2024, Wardana2024}.

\begin{figure}[H]
  \centering
  \fboxsep=8pt
  \fbox{\parbox{0.85\linewidth}{\centering\small
    \textbf{Gate-to-Gate System Boundary}\\[5pt]
    Washing $\;\rightarrow\;$ Soaking $\;\rightarrow\;$ Grinding $\;\rightarrow\;$ Boiling
    $\;\rightarrow\;$ Filtration $\;\rightarrow\;$ Coagulation $\;\rightarrow\;$ Molding \& Packaging
  }}
  \caption{Gate-to-gate system boundary of the tofu production life cycle assessment.}
  \label{fig:sysboundary}
\end{figure}

\subsection{Life Cycle Inventory (LCI) Analysis}

The inventory analysis phase involves the identification and collection of quantitative data on all
input and output flows for each unit process within the product system, in accordance with the defined
functional unit \citep{Saavedra2022}. Data collected include soybean inputs, water and energy
consumption, and waste outputs comprising liquid waste, solid waste, and atmospheric emissions. The
complete inventory dataset is presented in Table~\ref{tab:lci}.

\begin{table}[H]
\centering
\caption{Life Cycle Inventory Data of Tofu Production (per batch / production day).}
\label{tab:lci}
\small
\begin{tabularx}{\textwidth}{lXrclrc}
\toprule
\textbf{Stage} & \textbf{Input} & \textbf{Qty} & \textbf{Unit}
  & \textbf{Output} & \textbf{Qty} & \textbf{Unit}\\
\midrule
Washing       & Soybeans      & 1{,}400 & kg  & Washed soybeans   & 1{,}540 & kg  \\
              & Water         & 2{,}800 & L   & Wastewater        & 2{,}660 & L   \\
              & Electricity   & 0.60    & kWh & CO$_2$ emissions  & 0.51    & kg  \\
\midrule
Soaking       & Soybeans      & 1{,}540 & kg  & Hydrated soybeans & 3{,}640 & kg  \\
              & Water         & 2{,}700 & L   & Wastewater        & 600     & L   \\
              & Electricity   & 0.45    & kWh & CO$_2$ emissions  & 0.38    & kg  \\
\midrule
Grinding      & Hyd.\ soybeans& 3{,}640 & kg  & Soybean slurry    & 7{,}840 & kg  \\
              & Water         & 4{,}200 & L   & Wastewater        & 0       & L   \\
              & Electricity   & 17.00   & kWh & CO$_2$ emissions  & 14.45   & kg  \\
\midrule
Boiling       & Soybean slurry& 7{,}840 & kg  & Boiled slurry     &15{,}652 & kg  \\
              & Water         & 8{,}400 & L   & Steam loss        & 588     & L   \\
              & Firewood      & 800     & kg  & CO$_2$ emissions  & 1{,}397 & kg  \\
              &               &         &     & CH$_4$ emissions  & 0.0499  & kg  \\
              &               &         &     & N$_2$O emissions  & 0.3744  & kg  \\
\midrule
Filtration    & Boiled slurry &15{,}652 & kg  & Soy milk          &14{,}052 & kg  \\
              &               &         &     & Okara (residue)   & 1{,}600 & kg  \\
\midrule
Coagulation   & Soy milk      &14{,}052 & kg  & Tofu curd         & 3{,}681 & kg  \\
              & Coagulant water& 421    & L   & Whey              &10{,}434 & L   \\
\midrule
Molding \&    & Tofu curd     & 3{,}681 & kg  & Final product     & 3{,}681 & kg  \\
Packaging     & Water         & 100     & L   & Wastewater        & 100     & L   \\
\bottomrule
\end{tabularx}
\end{table}

\subsection{Life Cycle Impact Assessment (LCIA)}

The LCIA phase evaluates and quantifies the potential environmental impacts generated by the product
system based on the compiled inventory data \citep{Hauschild2017}. In this study, the
\textbf{CML-IA Baseline} method was applied using SimaPro~10.3, covering five impact categories:
Global Warming Potential (GWP$_{100a}$, kg CO$_2$~eq), Ozone Layer Depletion (ODP, kg CFC-11~eq),
Fresh Water Aquatic Ecotoxicity (FAET, kg 1,4-DB~eq), Acidification (AP, kg SO$_2$~eq), and
Eutrophication (EP, kg PO$_4^{3-}$~eq). Water use and potential water scarcity were additionally
assessed using the \textbf{AWARE (Available WAter REmaining)} method to evaluate water deprivation
relative to regional water availability.

\subsection{Interpretation}

The interpretation phase contextualizes LCIA results to identify hotspot impact categories and the
production stages contributing the greatest environmental burdens. Normalization was subsequently applied
using the CML-IA Baseline reference values to compare impact categories expressed in different units on
a common dimensionless scale, thereby enabling the relative significance of each impact to be
determined. Interpretation was focused strictly within the gate-to-gate system boundary to ensure that
results accurately reflect the stages exerting the greatest environmental influence \citep{Hauschild2017}.



% ══════════════════════════════════════════════════════════════════════════════
\section{Results and Discussion}
% ══════════════════════════════════════════════════════════════════════════════

\subsection{Environmental Impact of Tofu Production}

The LCIA reveals that tofu production generates measurable environmental impacts across all five
assessed categories. The quantitative results per production stage are summarized in
Table~\ref{tab:lcia}.

\begin{table}[H]
\centering
\caption{Environmental Impact of Tofu Production per Batch by Production Stage.}
\label{tab:lcia}
\small
\resizebox{\textwidth}{!}{%
\begin{tabular}{llllllllll}
\toprule
\textbf{Impact Category} & \textbf{Unit} & \textbf{Total}
  & \textbf{Washing} & \textbf{Soaking} & \textbf{Grinding}
  & \textbf{Boiling} & \textbf{Coagulation} & \textbf{Molding}\\
\midrule
GWP$_{100a}$   & kg CO$_2$~eq          & $5.28\!\times\!10^{3}$  & $3.56\!\times\!10^{3}$  & 0.945 & 35.8  & $1.63\!\times\!10^{3}$  & 55.8  & 5.9   \\
ODP            & kg CFC-11~eq          & $6.11\!\times\!10^{-5}$ & $5.93\!\times\!10^{-5}$ & $2.1\!\times\!10^{-9}$ & $7.93\!\times\!10^{-8}$ & $1.07\!\times\!10^{-6}$ & $5.02\!\times\!10^{-7}$ & $1.29\!\times\!10^{-7}$ \\
FAET           & kg 1,4-DB~eq          & $1.33\!\times\!10^{3}$  & $1.22\!\times\!10^{3}$  & 0.651 & 24.6  & 70.6                    & 14.1  & 2.75  \\
Acidification  & kg SO$_2$~eq          & 4.59   & 3.53    & 0.00245 & 0.0925 & 0.723 & 0.226 & 0.0176 \\
Eutrophication & kg PO$_4^{3-}$~eq     & 9.10   & 8.30    & 0.0034  & 0.128  & 0.409 & 0.250 & 0.00838 \\
\bottomrule
\end{tabular}}
\end{table}

\noindent\textit{Note: The filtration stage does not appear in the impact results because no direct
energy inputs or emissions were entered in the inventory model for this stage; it therefore generates
no characterization values in SimaPro.}

\subsubsection{Global Warming Potential (GWP)}

Total GWP amounts to \num{5.28e3}~kg CO$_2$~eq. The washing stage is the largest contributor at
\num{3.56e3}~kg CO$_2$~eq (approximately 67.4\%), followed by the boiling stage at
\num{1.63e3}~kg CO$_2$~eq (approximately 30.9\%). The dominance of the washing stage is attributed to
the large volumes of water consumed during soybean cleaning, which indirectly require energy for water
supply and distribution, thereby generating greenhouse gas emissions. In addition, organic-laden washing
wastewater has the potential to generate methane (CH$_4$) if discharged without treatment. The boiling
stage contributes significantly to GWP owing to the combustion of firewood, which releases CO$_2$,
CH$_4$, and N$_2$O, all of which are potent greenhouse gases with long-term atmospheric warming effects.

\subsubsection{Ozone Layer Depletion (ODP)}

Total ODP registers \num{6.11e-5}~kg CFC-11~eq, with the washing stage dominating at
\num{5.93e-5}~kg CFC-11~eq (approximately 97.1\%). This impact is not attributed to water consumption
per se, but rather to the energy-intensive infrastructure required for water supply and distribution
within the production chain, which generates indirect emissions linked to ozone-depleting substances.
Although this category contributes the least in normalized terms, its disproportionate concentration
in the washing stage highlights the indirect environmental burden of high water use.

\subsubsection{Fresh Water Aquatic Ecotoxicity (FAET)}

Total FAET reaches \num{1.33e3}~kg 1,4-DB~eq, with the washing stage as the primary contributor at
\num{1.22e3}~kg 1,4-DB~eq (approximately 91.7\%). The large volumes of organic-laden wastewater
generated during soybean washing contain residual organic matter, dissolved solids, and impurities that
constitute a significant toxicity risk to freshwater organisms if discharged without adequate treatment.
The boiling stage contributes 70.6~kg 1,4-DB~eq (approximately 5.3\%) through the atmospheric
deposition of combustion byproducts, including nitrogen oxides (NO$_x$) and sulfur oxides (SO$_x$),
into aquatic environments.

\subsubsection{Acidification}

Total Acidification impact is 4.59~kg SO$_2$~eq. The washing stage accounts for 3.53~kg SO$_2$~eq
(approximately 76.9\%), associated with SO$_2$ and NO$_x$ emissions from energy generation required
for water provision. The boiling stage contributes 0.723~kg SO$_2$~eq (approximately 15.8\%) through
firewood combustion, which releases acidifying gaseous precursors that may react with atmospheric
moisture to form acid deposition, potentially degrading both soil and water quality in the surrounding
environment.

\subsubsection{Eutrophication}

Total Eutrophication impact is 9.10~kg PO$_4^{3-}$~eq, with the washing stage dominant at
8.30~kg PO$_4^{3-}$~eq (approximately 91.2\%). This impact is driven primarily by nutrient-rich
washing wastewater that contains nitrogen and phosphorus compounds derived from soybean residues. When
discharged into receiving water bodies without treatment, these compounds can trigger excessive
proliferation of algae and aquatic microorganisms, leading to depletion of dissolved oxygen and
disruption of aquatic ecosystem function. The boiling stage contributes 0.409~kg PO$_4^{3-}$~eq
(approximately 4.5\%) through atmospheric deposition of NO$_x$ emissions generated during firewood
combustion.



\subsection{Water Use Assessment (AWARE Method)}

Water use was assessed using the AWARE method to quantify potential water scarcity, expressed as the
volume of water deprived from downstream users per unit of water consumed at a given location. The
total water scarcity impact for the entire production system amounts to \textbf{515~m$^3$ world~eq}
(Table~\ref{tab:aware}).

\begin{table}[H]
\centering
\caption{Water Use Assessment Results by Production Stage (AWARE Method).}
\label{tab:aware}
\small
\begin{tabular}{lcc}
\toprule
\textbf{Production Stage} & \textbf{Water Scarcity (m$^3$ world~eq)} & \textbf{Share (\%)}\\
\midrule
Boiling               & 294   & 57.09 \\
Grinding              & 101   & 19.61 \\
Soaking               & 39.3  & 7.63  \\
Coagulation           & 38.2  & 7.42  \\
Washing               & 35.8  & 6.95  \\
Molding \& Packaging  & 7.63  & 1.48  \\
\midrule
\textbf{Total}        & \textbf{515} & \textbf{100} \\
\bottomrule
\end{tabular}
\end{table}

The boiling stage is the dominant contributor (approximately 57.1\%), reflecting both direct water
consumption during soybean slurry cooking and the substantial upstream water demand embedded in the
firewood supply chain. Within the AWARE framework, the environmental burden is determined not solely by
the volume of water consumed, but also by the water intensity of the upstream energy system and
associated supply chains. The grinding stage contributes the second largest share (approximately
19.6\%), as continuous water addition is required throughout this process to achieve the appropriate
slurry consistency and extraction efficiency.

\subsection{Normalization}

Normalization was performed using the CML-IA Baseline reference values to place all impact categories
on a common dimensionless scale, enabling direct and meaningful comparison of relative environmental
significance across categories with different units (Table~\ref{tab:norm}).

\begin{table}[H]
\centering
\caption{Normalization Results of Tofu Production Environmental Impacts (CML-IA Baseline).}
\label{tab:norm}
\small
\begin{tabular}{lcc}
\toprule
\textbf{Impact Category} & \textbf{Normalized Value} & \textbf{Rank}\\
\midrule
Fresh Water Aquatic Ecotoxicity (FAET) & $2.57\times10^{-9}$  & \textbf{1} \\
Global Warming Potential (GWP)         & $1.05\times10^{-9}$  & 2 \\
Eutrophication                         & $6.89\times10^{-10}$ & 3 \\
Acidification                          & $1.63\times10^{-10}$ & 4 \\
Ozone Layer Depletion (ODP)            & $6.84\times10^{-13}$ & 5 \\
\bottomrule
\end{tabular}
\end{table}

Normalization confirms that \textbf{Fresh Water Aquatic Ecotoxicity (FAET) is the dominant
environmental impact category} in this production system, indicating that tofu production exerts the
greatest relative pressure on freshwater ecosystems. This finding underscores that organic liquid waste
management, particularly the treatment of effluents from the washing, soaking, and coagulation stages,
constitutes the most critical environmental challenge in the analyzed production system.

Although Global Warming Potential frequently receives primary attention in environmental sustainability
discourse, its normalized contribution ranks second in this study, indicating that freshwater pollution
represents a more acute environmental concern at the facility scale assessed. Eutrophication, which
ranks third in normalized impact, further reinforces the urgency of addressing nutrient-laden
wastewater streams as a priority intervention. Acidification and Ozone Layer Depletion, while present,
contribute relatively minor normalized values, suggesting that their management, though important, does
not constitute the primary environmental priority for this production system.



\subsection{Hotspot Identification and Improvement Strategies}

The integrated analysis of LCIA results and normalization outcomes identifies two primary environmental
hotspots in the tofu production system, namely the generation of organic liquid waste across the
washing, soaking, and coagulation stages, and the combustion of firewood during the boiling stage.
These hotspots correspond, respectively, to the dominant drivers of Fresh Water Aquatic Ecotoxicity
and Eutrophication (ranked first and third in normalized impact), and of Global Warming Potential and
water scarcity.

With respect to the first hotspot, the organic-laden wastewater streams generated during washing,
soaking, and coagulation are characterized by high BOD, COD, and nutrient loads
\citep{Seroja2018, Satar2022}, and their uncontrolled discharge poses significant risks to receiving
aquatic ecosystems. The implementation of anaerobic treatment systems, such as biodigesters or
integrated biogas wastewater treatment units, represents an effective and technically appropriate
intervention for SME-scale operations, as such systems simultaneously reduce organic loading while
generating biogas as a recoverable energy resource. Complementary measures include the application of
constructed wetland or biofilter systems for tertiary nutrient removal prior to effluent discharge, as
well as counter-current water recycling and recirculation within the washing and soaking stages to
reduce both volumetric water consumption and wastewater generation. Furthermore, the valorization of
whey, for instance as a protein-rich substrate for animal feed or as a feedstock for fermentation
processes, offers a viable pathway for diverting high-nutrient liquid streams away from the wastewater
treatment system, thereby reducing the overall nutrient load requiring treatment \citep{Bjornbet2021}.

With respect to the second hotspot, the boiling stage contributes disproportionately to GWP and water
scarcity owing to its reliance on firewood as the primary thermal energy source. The combustion of
800~kg of firewood per production batch releases 1{,}397~kg CO$_2$, 0.0499~kg CH$_4$, and
0.3744~kg N$_2$O, collectively representing a substantial greenhouse gas burden. Transitioning to
cleaner energy sources, such as liquefied petroleum gas (LPG) or biomethane derived from on-site
biodigesters, would directly reduce these emissions while also decreasing the upstream water footprint
associated with firewood supply chains. In parallel, the adoption of improved combustion systems or
thermally insulated boiling tanks would reduce specific fuel consumption per unit of output through
minimization of heat loss, while waste heat recovery from boiler exhaust could be directed toward
pre-heating process water at upstream stages, thereby improving overall energy efficiency. The
LCA-based analysis demonstrates that improvement interventions targeting liquid waste management and
energy use efficiency at the source generate substantially greater environmental benefits than
conventional end-of-pipe approaches \citep{Bjornbet2021, Apriyanti2025}.

% ══════════════════════════════════════════════════════════════════════════════
\section{Conclusion}
% ══════════════════════════════════════════════════════════════════════════════

This study applied Life Cycle Assessment with a gate-to-gate system boundary to evaluate the
environmental performance of small-scale tofu production at an SME in Semarang City. The assessment
encompassed five impact categories under the CML-IA Baseline method, supplemented by water scarcity
analysis using the AWARE method, and yielded quantitative evidence on the distribution of environmental
burdens across all production stages.

The results demonstrate that tofu production generates measurable impacts across all assessed
categories, with total values of \num{5.28e3}~kg CO$_2$~eq for GWP, \num{6.11e-5}~kg CFC-11~eq for
ODP, \num{1.33e3}~kg 1,4-DB~eq for FAET, 4.59~kg SO$_2$~eq for Acidification, and
9.10~kg PO$_4^{3-}$~eq for Eutrophication. The washing and boiling stages consistently emerge as the
principal contributors across all impact categories, driven by high water consumption and firewood
combustion, respectively. Normalization of these results reveals that Fresh Water Aquatic Ecotoxicity
constitutes the dominant environmental impact category, with a normalized value of
$2.57\times10^{-9}$, confirming that the organic liquid waste streams generated throughout the
production process represent the most critical environmental hotspot. Water scarcity, as assessed by
the AWARE method, totals 515~m$^3$ world~eq, with the boiling stage accounting for 57.1\% of this
burden, a value that reflects not only direct water use but also the embedded water demand of firewood
procurement and combustion.

These findings indicate that environmental improvement efforts should be strategically directed toward
two priority areas. The first is the establishment of structured liquid waste management systems,
including anaerobic treatment, nutrient recovery, and water recycling, to address the dominant
freshwater ecotoxicity and eutrophication burdens arising from washing, soaking, and coagulation
effluents. The second is the substitution of firewood with cleaner fuel alternatives and the adoption
of energy-efficient thermal systems in the boiling stage, which would simultaneously reduce greenhouse
gas emissions, acidification potential, and water scarcity impacts. These conclusions are consistent
with broader findings in the LCA literature on soy-based food processing, which collectively emphasize
the critical importance of wastewater management and energy source selection as the primary levers for
environmental performance improvement in tofu production systems \citep{Nugroho2024, Wardana2024}.

The findings of this study provide actionable, evidence-based guidance for plant operators and
policymakers seeking to improve the environmental sustainability of SME-scale tofu production, and may
inform the development of sector-specific environmental management standards for the tofu industry in
Semarang City and Central Java.



% ── Bibliography ──────────────────────────────────────────────────────────────
\noindent\textcolor{ruleblue}{\rule{\linewidth}{0.6pt}}

\begin{small}
\noindent\textbf{Reference Validation Notes:} Entries marked \textcolor{warnamber}{$\star$} could not
be independently verified and should be confirmed prior to journal submission.
\end{small}

\vspace{4pt}

\begin{thebibliography}{15}
\small

\bibitem{Ghisellini2016}
Ghisellini, P., Cialani, C., \& Ulgiati, S. (2016).
A review on circular economy: The expected transition to a balanced interplay of environmental and
economic systems.
\textit{Journal of Cleaner Production}, \textit{114}, 11--32.
\url{https://doi.org/10.1016/j.jclepro.2015.09.007}

\bibitem{Dewi2024}
\textcolor{warnamber}{$\star$}~Dewi, A.~N., \& Setiawan, D. (2024).
Bisnis kuliner: Studi kasus CV.\ Gehu Extra Pedas Chili Hot.
\textit{Jurnal Manajemen dan Bisnis Islam}, \textit{1}(1), 5--23.
[\textit{Note: Cited for BPS per capita consumption data; journal relevance to LCA data source could
not be independently verified. Consider replacing with a direct BPS statistical publication.}]

\bibitem{Rosyidah2020}
Rosyidah, M., Masruri, A., \& Putra, R.~A. (2020).
Assessment (LCA) method on tofu production.
\textit{International Journal of Science, Technology and Management}, \textit{1}(4), 428--435.

\bibitem{Sjafruddin2022}
Sjafruddin, R., Agustang, A., \& Pertiwi, N. (2022).
Estimasi limbah industri tahu dan kajian penerapan sistem produksi bersih.
\textit{Jurnal Ilmiah Mandala Education}, \textit{8}(2), 1229--1237.
\url{https://doi.org/10.36312/jime.v8i2.2826}

\bibitem{Islami2025}
\textcolor{warnamber}{$\star$}~Islami, M.~C.~P.~A., \& Harnaningrum, R.~N. (2025).
Integration of waste management and environmental impact assessment for sustainable manufacturing in
Sukolego tofu production.
\textit{[Journal name not specified in source document]}, \textit{6}(2), 71--77.
[\textit{Note: Journal name is missing. Authors must complete this reference before submission.}]

\bibitem{Hauschild2017}
Hauschild, M.~Z., Rosenbaum, R.~K., \& Olsen, S.~I. (2017).
\textit{Life Cycle Assessment: Theory and Practice}.
Springer, Cham.
\url{https://doi.org/10.1007/978-3-319-56475-3}

\bibitem{Chitaka2023}
Chitaka, T.~Y., \& Goga, T. (2023).
The evolution of life cycle assessment in the food and beverage industry: A review.
\textit{Cambridge Prisms: Plastics}, \textit{1}, 1--6.
\url{https://doi.org/10.1017/plc.2023.4}

\bibitem{Nugroho2024}
Nugroho, M.~E., Setyono, P., \& Rachmawati, S. (2024).
Analisis emisi gas rumah kaca dengan Life Cycle Assessment (LCA) dan Analytical Hierarchy Process
(AHP) industri tahu.
\textit{Jurnal Ilmu Lingkungan}, \textit{22}(6), 1504--1512.
\url{https://doi.org/10.14710/jil.22.6.1504-1512}

\bibitem{Wardana2024}
Kartika Wardana, S., Cucikodana, Y., Almaniar, S., Dwijayanti, A., Maulana, F., \&
Muhlisoh, N.~A. (2024).
Penilaian dampak lingkungan dengan Life Cycle Assessment (LCA) pada industri tahu Kampung Jangkar
Kulon, Cilegon Banten.
\textit{Jurnal Teknologi Kimia Unimal}, \textit{13}(2), 97--106.
\url{https://doi.org/10.29103/jtku.v13i2.16429}

\bibitem{Bjornbet2021}
Bj{\o}rnbet, M.~M., \& Vild{\aa}sen, S.~S. (2021).
Life cycle assessment to ensure sustainability of circular business models in manufacturing.
\textit{Sustainability}, \textit{13}(19), article~11014.
\url{https://doi.org/10.3390/su131911014}

\bibitem{Saavedra2022}
Saavedra-Rubio, K., Thonemann, N., Crenna, E., Lemoine, B., Caliandro, P., \& Laurent, A. (2022).
Stepwise guidance for data collection in the life cycle inventory (LCI) phase: Building
technology-related LCI blocks.
\textit{Journal of Cleaner Production}, \textit{366}, article~132903.
\url{https://doi.org/10.1016/j.jclepro.2022.132903}

\bibitem{Rasyid2024}
\textcolor{warnamber}{$\star$}~Rasyid, M., \& Anggriani, R. (2024).
Penerapan Life Cycle Assessment (LCA) pada proses produksi minyak kayu putih di Desa Sawa-Namlea.
\textit{Journal of Social Science Research}, \textit{4}(3), 18970--18984.
[\textit{Note: Journal peer-review status and indexing could not be independently confirmed. Consider
replacing with ISO 14044:2006 or a verified LCA methodology reference.}]

\bibitem{Seroja2018}
Seroja, R., Effendi, H., \& Hariyadi, S. (2018).
Tofu wastewater treatment using vetiver grass (\textit{Vetiveria zizanioides}) and zeliac.
\textit{Applied Water Science}, \textit{8}(1), 1--6.
\url{https://doi.org/10.1007/s13201-018-0640-y}

\bibitem{Satar2022}
Satar, I., \& Permadi, A. (2022).
Treating the tofu wastewater (TWW) using a green technology of microbial fuel cell (MFC) system.
\textit{Indonesian Journal of Environmental Management and Sustainability}, \textit{6}(1), 162--167.

\bibitem{Apriyanti2025}
Apriyanti, D., Pratikno, F.~A., \& Hertadi, C.~D.~P. (2025).
Penilaian dampak lingkungan menggunakan Life Cycle Assessment (LCA) pada proses produksi tempe di SIKS.
\textit{Jurnal Mitra Teknik Industri}, \textit{4}(1), 19--27.
\url{https://doi.org/10.24912/jmti.v4i1.34600}

\end{thebibliography}

\end{document}
